% ============================================================
%  ch11  第四讲 §4.3 adèle 单值化
%  原文行 2293–2457
% ============================================================

\subsection{adèle 单值化}

我们已看到用 $\Gr_x$ 单值化 $\Bun_{\underline G}$ 的一些重要应用。但在实践中，该单值化定理有几处局限。首先它只对半单群有效；更严重的是，$u_x$ 的纤维（它们是 $\underline G^{X^*}$-主丛）的几何仍然复杂，因此关于 $\Bun_{\underline G}$ 的某些问题难以通过研究 $\Gr_x$ 得到（如定理 \ref{uniformize line}）。本小节引入 $\Bun_{\underline G}$ 的另一种单值化来克服这些困难。

回忆 A. Weil 的下述结果。设 $k$ 代数闭，$X$ 是 $k$ 上光滑连通射影曲线，$|X|$ 记闭点集。设 $\bO=\prod_{x\in |X|} \mO_x$ 是整 adèle 环，$\bA=\prod' F_x$ 是 adèle 环，即 $F_x$（$x\in |X|$）关于所有闭点的\emph{限制}乘积：$\bA$ 中的元素是族 $(f_x\in F_x)_{x\in |X|}$，其中除有限多个 $x$ 外均有 $f_x\in \mO_x$。设 $\eta$ 是 $X$ 的通用点，$F_\eta$ 是 $X$ 的函数域。注意自然嵌入 $F_\eta\subset F_x$ 给出对角嵌入 $F_\eta\subset \bA$。

设 $\underline G$ 是 $X$ 上逐纤维连通的平滑仿射群概形，记 $G=\underline G\otimes F_\eta$。则

\begin{lem}
有自然双射
\begin{equation}\label{Weil}
\Bun_{\underline G}(k)\simeq G(F_\eta)\backslash G(\bA)/\underline G(\bO).
\end{equation}
\end{lem}

\begin{proof}
设 $\mE$ 是 $X$ 上的 $\underline G$-主丛。由于 $F_\eta$ 是上同调维数为一的域，由 Steinberg 定理，$G$-主丛 $\mE\otimes F_\eta$ 平凡。取这样的平凡化 $\beta_\eta: \mE\otimes F_\eta\simeq \mE^0\otimes F_\eta$，它可延拓为某个开子集 $U\subset X$ 上的平凡化。

另一方面，对每个 $x\in |X|$，$\mE|_{D_x}$ 也平凡，固定这样的平凡化 $\beta_x: \mE|_{D_x}\simeq \mE^0|_{D_x}$。则 $(\beta|_{D_x^*})\circ(\beta_x^{-1}|_{D_x^*})$ 是 $D_x^*$ 上平凡 $\underline G$-主丛的自同构，从而定义一个元素 $g_x\in G(F_x)$。注意对 $x\in |U|$ 有 $g_x\in G(\mO_x)$，于是由三元组 $(\mE,\beta_\eta,\{\beta_x\}_{x\in |X|})$ 构造出元素 $(g_x)\in G(\bA)$。$\beta_\eta$ 与 $\beta_x$ 的不同选取，会把该元素分别左乘 $G(F_\eta)$ 中元素、右乘 $G(\bO)$ 中元素加以改变。因此 $\mE$ 给出 $G(F_\eta)\backslash G(\bA)/\underline G(\bO)$ 中一个确定的元素，这就定义了映射 $\Bun_{\underline G}(k)\to G(F_\eta)\backslash G(\bA)/\underline G(\bO)$。

反之，给定 $g=(g_x)\in G(\bA)$，存在开子集 $U\subset X$ 使 $x\in |U|$ 时 $g_x\in G(\mO_x)$。于是可用 $g_x\in G(F_x)=\Aut(\mE^0|_{D_x^*})$ 作为过渡函数，把 $U$ 上与 $D_x$（$x\in |X|\setminus |U|$）上的平凡主丛粘起来，得到 $X$ 上的 $\mE$。该构造给出逆映射。
\end{proof}

下面描述 \eqref{Weil} 的代数几何版本，称之为 $\Bun_{\underline G}$ 的\textbf{adèle 单值化}。第一个问题是如何赋予集合 $G(\bA)/\underline G(\bO)$ 代数几何结构。事实上，这个问题有两种自然的解决方案。

第一种方案用 Ran Grassmann 簇 $\Gr_{\Ran}$（定义 \ref{dfn Ran Grass}）。$G(\bA)/\underline G(\bO)$ 中的元素可解释为一对 $(\mE,\beta_\eta)$，其中 $\mE$ 是 $X$ 上的 $\underline G$-主丛，$\beta_\eta$ 是 $\mE$ 在 $\eta$ 上（等价地，在某个开子集 $U\subset X$ 上）的截面；而 $\Gr_{\Ran}$ 的 $k$-点还指定了 $X$ 上的一组点，$\beta_\eta$ 在这组点上可能不再是正则映射。可按如下方式去掉这种冗余（沿 \cite[\S~2]{Ga3} 的做法）。

回忆 Ran 空间 $\Ran=\Ran(X)$ 有点集之并给出的半群结构。有「作用」映射
\[\on{act}: \Ran\times \Gr_{\Ran}\to \Gr_{\Ran},\]
把 $(x, (y,\mE,\beta))\in \Ran(R)\times\Gr_{\Ran}(R)$ 送到 $(x\cup y, \mE, \beta|_{X_R-\Ga_{x\cup y}})$。所谓「商」预层记作 $\Gr_{\Ran,\on{indep}}$，定义为下图的上等化子
\[\Ran\times\Gr_{\Ran}\substack{\on{act}\\  \xrightarrow{\hspace*{.6cm}} \\[-0.6ex]  \xrightarrow{\hspace*{.6cm}} \\ \pr_2 } \Gr_{\Ran}.\]
具体说来，$\Gr_{\Ran,\on{indep}}(R)$ 是 $\Gr_{\Ran}(R)$ 按如下等价关系的商：$(x,\mE,\beta)\sim (x',\mE',\beta')$，若存在 $X_R$ 上 $\underline G$-主丛的同构 $\al: \mE\cong \mE'$，使限制到 $X_R-\Ga_{x\cup x'}$ 时与 $\beta$、$\beta'$ 相容。由定义即得
\[\Gr_{\Ran,\on{indep}}(k)= G(\bA)/\underline G(\bO).\]
但一般地 $\Gr_{\Ran,\on{indep}}$ 不是空间，除非 $\underline G$ 是平凡群（此时 $\Gr_{\Ran,\on{indep}}$ 是一点）。所幸 $\Gr_{\Ran,\on{indep}}$ 的 fpqc 层化存在，并可显式描述如下。

称开子概形 $U\subset X_R$ \textbf{极稠密}，若它逐纤维稠密，即复合 $U\to X_R\to \Spec R$ 忠实平坦。注意极稠密开子集的有限交仍极稠密。

定义 $\Gr_{\underline G,\Rat(X)}$ 为给每个 $k$-代数 $R$ 配上下述群胚的预叠：对象是三元组 $(\mE,U,\beta)$，其中 $\mE$ 是 $X_R$ 上的 $\underline G$-主丛，$U\subset X_R$ 是极稠密开子概形，$\beta$ 是 $\mE$ 在 $U$ 上的平凡化；从 $(\mE,U,\beta)$ 到 $(\mE',U',\beta')$ 的态射是使 $\beta'\alpha|_{U\cap U'}=\beta|_{U\cap U'}$ 的同构 $\alpha:\mE\cong \mE'$。

若 $\underline G$ 与 $X$ 在上下文中明确，为简洁把 $\Gr_{\underline G,\Rat(X)}$ 记作 $\Gr_{\Rat}$。需要

\begin{lem}\label{rat sheaf}
上述群胚离散，且 $\Gr_{\Rat}$ 实际上是层。
\end{lem}

\begin{proof}
要证任何对象都没有非平凡自同构，即：设 $R$ 是 $k$-代数，$\mE$ 是 $X_R$ 上的 $\underline G$-主丛，$\al: \mE\to \mE$ 是自同构，若 $\al|_U=\id$，则 $\al=\id$。把 $\mE$ 嵌入 $X_R$ 上某个仿射空间，便化归为：若 $f$ 是 $X_R$ 上的正则函数，且在某个极稠密开子集上为零，则 $f=0$。注意这并不完全平凡，因为没有假设 $R$ 既约。它由下述引理推出。
\end{proof}

\begin{lem}\label{very dense vs Cartier}
设 $X$ 是拟射影簇，$Z\subset X$ 是有限（可为空）闭点集，看作闭子概形。设 $U\subset X_R$ 是开稠密子集，满足 (i) $Z_R\subset U$；(ii) $U\to \Spec R$ 满射。则存在 $\Spec R$ 由 $\Spec R_i$ 组成的开覆盖与相对 Cartier 除子 $D_i\subset X_{R_i}$，使 $Z_R\cap D_i=\emptyset$ 且 $X_{R_i}\setminus D_i\subset U\otimes_RR_i$。
\end{lem}

回忆在仿射空间 $X=\Spec A$ 上，Cartier 除子按定义是由某个\emph{非}零因子 $f\in A$ 的方程 $f=0$ 定义的闭子概形（参见 \cite[Tag01WQ]{St}）。于是由定义，局部化映射 $A\to A_f$ 单射，因此在 $X\setminus D$ 上为零的函数在整个 $X$ 上为零。

\begin{proof}
化归到如下情形：$X=\Spec A$ 仿射，$S=\Spec R$，$U=\Spec A\otimes_k R\setminus V(\fraka)$，其中 $\fraka$ 是 $A\otimes_k R$ 的理想，且 $V(\fraka)\cap Z=\emptyset$。

$U\to S$ 的满射性等价于：对 $R$ 的每个素理想 $\frakp$，令 $K=\on{Frac}(R/\frakp)$，有 $\on{Im}(\fraka\otimes_RK\to A\otimes_kK)\neq 0$。固定这样的 $\frakp$，并可自由缩小 $S$ 使其围绕 $x=V(\frakp)$。

把 $R$ 写成有限生成 $k$-代数之并，并缩小 $U$，可设 $R$ 有限生成且 $U\to \Spec R$ 的像含 $x$。取 $\on{Im}(\fraka\otimes_R R/\frakp\to A\otimes_k R/\frakp)$ 中这样的非零元 $\bar{f}$，使 $V(\bar{f})\cap Z_K=\emptyset$，并取 $\bar f$ 的提升 $f\in \fraka$。则由 \cite[Tag00ME]{St}，乘法映射 $f: A\otimes_k R_\frakp\to A\otimes_k R_\frakp$ 单射且商为 $R_\frakp$-平坦。因 $R$ noether，必要时缩小 $S$ 使其围绕 $x$，可设 $R\to R_\frakp$ 单射，于是 $f:A\otimes_k R\to A\otimes_k R$ 单射。再次必要时缩小 $S$，可设 $V(f)$ 不含竖直除子且 $V(f)\cap Z=\emptyset$。于是对 $R$ 的每个素理想 $\frakq$，$f: A\otimes_k R/\frakq\to A\otimes_k R/\frakq$ 单射。因此 $(A\otimes_k R/f)_\frakq$ 对每个 $R_\frakq$ 平坦，从而 $A\otimes_kR/f$ 在 $R$ 上平坦。则 $D=V(f)$ 即所求的相对 Cartier 除子。
\end{proof}

显然，映射 $(x,\mE,\beta)\mapsto (\mE,U=X_R-\Ga_x,\beta)$ 把 $\Gr_{\Ran,\on{indep}}$ 实现为 $\Gr_{\on{Rat}}$ 的子预层。

\begin{lem}\label{sheafify Ran}
空间 $\Gr_{\on{Rat}}$ 是 $\Gr_{\Ran,\on{indep}}$ 的 fpqc 层化。
\end{lem}

\begin{proof}
只需证明 fpqc 局部地，每个点 $\Spec R\to \Gr_{\on{Rat}}$ 都能提升为 $\Gr_{\Ran,\on{indep}}$ 的点。设 $(\mE,U, \beta)$ 代表 $\Gr_{\on{Rat}}$ 的 $R$-点。由引理 \ref{very dense vs Cartier}，在 $\Spec R$ 的 Zariski 局部可设 $U=X_R\setminus D$，其中 $D\subset X_R$ 是相对 Cartier 除子。回忆给每个 $R$ 配上 $X$ 上相对 Cartier 除子集的预层由 $X$ 的对称幂表示。另一方面，对称化映射 $X^d\to X^{(d)}$ 有限忠实平坦。因此经 $R$ 的有限忠实平坦基变换后可设 $U$ 是某个 $x:\Spec R\to X^d$ 的 $\Ga_x$ 的补集。
\end{proof}

其次需要 $G(F_\eta)$ 上的代数几何结构，同样有两种方案。这里只介绍 ${\underline G}_{\Rat}:={\underline G}_{\Rat(X)}$，即对应于 $\Gr_{\on{Rat}}$ 的那种。它给每个 $R$ 配上对 $(U,\al)$ 的群胚，其中 $U$ 是 $X_R$ 上的极稠密开子集，$\al\in \underline G(U)$。这仍是群层。下述引理容易。

\begin{lem}\label{generic dependence}
${\underline G}_{\Rat(X)}$ 只依赖 $\underline G\otimes F_\eta$，且 ${\underline G}_{\Rat(X)}(k)=G(F_\eta)$。
\end{lem}

现在可陈述 \eqref{Weil} 的代数几何版本。

\begin{thm}\label{ran uniformization of BunG}
设 $\underline G$ 的通用纤维约化。则态射 $u_{\Rat}:\Gr_{\Rat}\to\Bun_{\underline G}$ 是 ${\underline G}_{\Rat}$-主丛，从而诱导同构
\[[{\underline G}_{\Rat(X)}\backslash \Gr_{\Rat}]\cong \Bun_{\underline G}.\]
确切地说，对每个 $\Spec R\to \Bun_{\underline G}$，存在忠实平坦映射 $\Spec R'\to \Spec R$，使 $\Spec R'\times_{\Bun_{\underline G}}\Gr_{\Rat}$ 与 $\Spec R'\times {\underline G}_{\Rat}$ 在 ${\underline G}_{\Rat}$-等变意义下同构。
\end{thm}

\begin{proof}
基本要素与定理 \ref{uniformization for P1} 相同：$X_R$ 上每个 $\underline G$-主丛在忠实平坦基变换 $\Spec R'\to\Spec R$ 后，都在某个极稠密开子集上可约化为 Borel。为证此，由引理 \ref{generic dependence} 可修改 $\underline G$ 的某些纤维使其成为 $X$ 上的 Bruhat--Tits 群概形（定义见注 \ref{projectivity of BD}），再应用 \cite[Corollary 26]{He}。又因 $\bG_m$-与 $\bG_a$-主丛可 Zariski 局部平凡，定理得证。
\end{proof}

引入上述 $\Bun_{\underline G}$ 的 adèle 单值化定理，一个原因是 $u_{\Rat}$ 的纤维的几何出人意料地「简单」：群空间 ${\underline G}_{\Rat(X)}$ 在某种意义上表现得像一点。下面用定理来支持这一断言。回忆对 $k$-空间 $\mF$（或 $k$-预叠），可以谈论其正则函数环与线丛的 Picard 群胚（见 \eqref{Pic line}）；\cite[\S~2]{GL} 还解释了如何定义预叠的（$\ell$-进）（上）同调。

\begin{thm}
设 $R$ 是有限型 $k$-代数。
\begin{enumerate}
\item[(1)] $\underline G_{\Rat}\otimes R$ 上每个正则函数都是 $\Spec R$ 上某正则函数的拉回。
\item[(2)] $\underline G_{\Rat}\otimes R$ 上每个线丛都是 $\Spec R$ 上某线丛的拉回。
\item[(3)] 映射 $\underline G_{\Rat}\otimes R\to \Spec R$ 诱导（$\ell$-进）同调的同构。
\end{enumerate}
\end{thm}

(1) 与 (2) 归功于 Beilinson--Drinfeld（见 \cite[Proposition 4.3.13]{BD} 与 \cite[\S~4.9.1]{BD0}）；(3) 归功于 Gaitsgory--Lurie（见 \cite{Ga3} 与 \cite[Theorem 3.3.1]{GL}）。

\begin{proof}
(3) 参见 \cite{GL}；这里仅对 $G=\bG_m$ 或 $\bG_a$ 的情形解释 (1)，一般情形化归到该特殊情形（用引理 \ref{generic dependence}）。

取 $X$ 上的丰沛线丛 $\mL$。对 $n\gg 0$，令 $V_n=\Gamma(X,\mL^n)$，$V'_n= V_n\setminus \{0\}$。则有映射
\[p_n:V'_n\times V'_n\to (\bG_m)_{\on{rat}(X)},\quad  p_n:V_n\times V'_n\to (\bA^1)_{\on{rat}(X)},\]
由 $p_n(f,g)\mapsto f/g$ 给出。由引理 \ref{very dense vs Cartier}，每个映射 $S\to (\bG_m)_{\on{rat}(X)}$（或 $S\to (\bA^1)_{\on{rat}(X)}$）在 $S$ 的 Zariski 局部，对某个足够大的 $n$ 通过 $p_n$ 分解。

给定 $Y_{\on{rat}(X)}\otimes R$ 上的正则函数 $\varphi$，其拉回给出 $(V'_n\times V'_n)\otimes R$（或 $(V_n\times V'_n)\otimes R$）上的正则函数 $p_n^*\varphi$，它关于 $\bG_m$ 在 $V_n'\times V'_n$（或 $V_n\times V'_n$）上的显然作用不变。取 $n$ 足够大使 $\dim V_n>1$。则 $p_n^*\varphi$ 延拓为 $(V_n\times V_n)\otimes R$ 上 $\bG_m$-不变的正则函数，而它必来自 $R$。这就证明了 (1)。(2) 可类似论证。
\end{proof}

为使该定理更有用，需要下述配套陈述。按定义有映射序列
\[\Gr_{\Ran}\stackrel{w}{\to} \Gr_{\Rat}\stackrel{u_{\on{Rat}}}{\to} \Bun_{\underline G},\]
且 $u_{\Ran}=w\cdot u_{\on{Rat}}$。有

\begin{thm}\label{Ran=Rat}~
\begin{enumerate}
\item[(1)] 沿 $w$ 的拉回诱导正则函数空间的同构。
\item[(2)] 沿 $w$ 的拉回诱导刚性化线丛的 Picard 群胚的等价。
\item[(3)] 沿 $w$ 的推前诱导 $\ell$-进同调的同构。
\end{enumerate}
\end{thm}

注意即使 $\underline G$ 是平凡群，该定理也不平凡；此时把它陈述为下述命题。

\begin{prop}\label{Ran pt}~
\begin{enumerate}
\item[(1)] $\Ran X\otimes R$ 上每个正则函数都是 $\Spec R$ 上某正则函数的拉回。
\item[(2)] 自然映射 $\Ran X\otimes R\to \Spec R$ 诱导（$\ell$-进）同调的同构。
\end{enumerate}
\end{prop}

\begin{proof}
(2) 仍参见 \cite{Ga3} 与 \cite[\S~2]{GL}。为证 (1)，用下述引理。
\begin{lem}
设 $f\in R[[x,y,z]]$ 是系数在 $R$ 中的关于 $x,y,z$ 的幂级数。若 $f$ 关于 $x,y,z$ 的置换不变，且 $f(x,x,y)=f(x,y,y)$，则 $f(x,x,x)$ 为常数。
\end{lem}
该引理留作有趣的习题。把 $X^I\to \Ran(X)$ 分解为 $X^I\stackrel{\Delta}{\to} X^{I\sqcup I\sqcup I}\to \Ran(X)$，便知 $\Ran\otimes R$ 上每个函数拉回到 $X^I\otimes R$ 时都来自 $\Spec R$ 上的正则函数，于是 (1) 成立。
\end{proof}

下面解释定理 \ref{Ran=Rat} 的 (1)。回忆预叠上正则函数空间的定义（参见 \eqref{reg fun}）。由于 $\bG_a$ 是层，由引理 \ref{sheafify Ran}，
\[\Gamma(\Gr_{\Rat},\mO)=\Hom(\Gr_{\Rat},\bG_a)=\Hom(\Gr_{\Ran,\on{indep}},\bG_a)=\Gamma(\Gr_{\Ran,\on{indep}},\mO),\]
而按定义后者同构于下述序列的等化子
\[\Gamma(\Gr_{\Ran},\mO)\substack{\on{act}^*\\  \xrightarrow{\hspace*{.6cm}}\\[-0.6ex] \xrightarrow{\hspace*{.6cm}} \\ \pr_2^*} \Gamma(\Ran\times \Gr_{\Ran},\mO).\]
因此要证 (1)，只需说明 $\on{act}^*=\pr_2^*$。有自然映射 $(q_{\Ran},\id):\Gr_{\Ran}\to \Ran\times\Gr_{\Ran}$，把 $(x,\mE,\beta)$ 送到 $(x,(x,\mE,\beta))$；此外
\[ \on{act}\circ(q_{\Ran},\id)=\pr_2\circ(q_{\Ran},\id)=\id.\]
由命题 \ref{Ran pt} (1)，$\pr_2^*$ 是同构，故 $(q_{\Ran},\id)^*$ 是同构，从而 $\on{act}^*=\pr_2^*$。定理 \ref{Ran=Rat} 的 (2) 可类似证明。

\begin{cor}\label{Gr=Bun}~
\begin{enumerate}
\item 沿 $u_{\Ran}$ 的拉回诱导正则函数空间的同构。
\item 它诱导刚性化线丛的 Picard 群胚的等价。
\item 推前诱导 $\Gr_{\Ran}$ 与 $\Bun_{\underline G}$ 的同调的同构。
\end{enumerate}
\end{cor}

\begin{rmk}
(i) 注意把定理 \ref{Rigid line} 与推论 \ref{Gr=Bun} 的 (2) 结合，可得定理 \ref{uniformize line} 的自然证明。

(ii) 推论 \ref{Gr=Bun} 的 (3) 是 Gaitsgory--Lurie 证明整体函数域上 Weil 关于 Tamagawa 数猜想的关键步骤之一（参见 \cite{GL}）。
\end{rmk}
